\chapter{Review of Related Literature}
\label{ch:literature}

\section{Process representation determines the meaning of a prediction}

\subsection{Treatment purpose and the scope of engineering evidence}

Wastewater treatment must address pollutants with different sources, transformations, and effects. \citet{Madhav2019} describe these pollutant sources and environmental consequences, while \citet{Lin2022} examine differences in their health effects. Their accounts explain why one effluent indicator cannot represent every treatment concern. An accurate COD prediction, for example, leaves the nitrogen distribution and retained solids unresolved. Prediction quality therefore needs to be judged against the treatment question that the output is intended to answer.

The engineering purpose is also shaped by implementation conditions. \citet{Tabios2020} relates Philippine water resources to physical conditions and governance. \citet{DamalerioBarriers2022} identify technical, financial, social, and institutional barriers to decentralized wastewater management. The practical guide edited by \citet{Ulrich2009} connects treatment selection to operation and maintenance in developing countries. These sources broaden the meaning of a useful treatment decision. Faster numerical evaluation can support technical comparison, but its speed is not evidence that financing, staffing, or institutional barriers have been resolved. This distinction places computational advances within their appropriate contribution to engineering practice.

Treatment frameworks consequently consider several priorities. \citet{Jafarinejad2020} places energy efficiency, cost, reliability, resilience, and sustainability within a common design discussion. \citet{Narayanan2019} connect biological processes to bioreactor design. \citet{Gu2023} review the competing demands involved in aeration optimization and control. These contributions make a single performance ranking inadequate without its stated objective. A decision that reduces oxygen-transfer capacity may carry a different nutrient response. Evidence about a model should establish which of these quantities it predicts and which trade-offs its assessment actually examines.

\subsection{From sludge inventories to component-based nutrient models}

The development of activated sludge models provides a basis for defining the predicted state. \citet{Marais1976} organize steady-state behavior around substrate, biomass, and operating conditions. \citet{Dold1981} develop a more general process representation. \citet{Henze1987} establish the single-sludge framework associated with Activated Sludge Model No.\ 1 (ASM1). These contributions connect reported treatment quantities to internal biological and chemical states. Their distinction remains important for learning. A measured total is an observation of a state, while the complete state contains several quantities whose changes must remain mutually consistent.

Different models introduce different internal explanations for a similar external response. \citet{Gujer1999} develop Activated Sludge Model No.\ 3 (ASM3) with explicit intracellular storage. \citet{Rieger2001} extend that storage-based representation to biological phosphorus removal. \citet{Fall2015} modify an activated sludge model to represent low solids production. These developments improve the relevance of a model to a particular process question, but they also change the quantities whose mass conservation must be checked. An added storage pool or altered yield cannot be treated as a new label attached to an unchanged state representation.

Nutrient-removal formulations make this dependence especially clear. \citet{Barker1997} present a general biological nutrient-removal model, and \citet{Hu2007} develop its kinetic representation. \citet{Henze1999} describe a phosphorus-removal model that includes denitrification by phosphorus-accumulating organisms. The collected models in \citet{Henze2006} and the treatment principles in \citet{Henze2008} show how these choices relate to engineering conditions. Their differences concern the represented transformations as well as the fitted parameters. A surrogate comparison that changes the component basis between learners also changes the prediction problem and weakens an attribution of differences to the learning algorithms alone.

The long-term assessment of ASM1 by \citet{VanLoosdrecht2015} and the critical review by \citet{Hauduc2013} discourage equating additional model detail with greater credibility. More detail can describe additional mechanisms, but it also increases the information needed for calibration. A simpler model can be adequate for a defined output without representing every internal pathway. The defensible comparison is therefore between the scientific question, the represented state, and the evidence available to identify that state. The model's name does not settle this comparison.

Constituent accounting is another representation choice. \citet{Hoover1952} relate biomass synthesis to oxygen utilization and the composition of assimilated material. \citet{TakacsVanrolleghem2006} examine elemental balances in activated sludge modeling, including the consequences of omitted products. These contributions expose a limitation in loosely stated mass conservation claims. The quantities represented by mass conservation are defined by the adopted lumped-component basis. Accounting for them does not establish complete accounting for every unrepresented element or external stream. The reporting boundary and the boundary for mass conservation must be identified before a physical guarantee can be interpreted.

\subsection{Biological pathways and transport place limits on transfer}

Carbon and nitrogen removal are linked through microbial metabolism and substrate availability. \citet{Wentzel1992} explain the connections among nitrification, denitrification, and excess phosphorus removal. \citet{Guerrero2011} examine competition between phosphorus-accumulating organisms and denitrifiers under different carbon sources. \citet{Lu2014} review the microbial ecology of denitrification, while \citet{Duan2022} examine heterotrophic nitrification. Together, these sources show why equal reported totals can conceal different process states. The identity of the carbon source and the form of nitrogen affect the reactions that remain possible. Independent predictions of aggregated quantities can consequently appear accurate while omitting relationships needed to interpret an operating change.

Functional process descriptions also require care. \citet{Preena2021} review nitrification and denitrification in recirculating aquaculture, which supplies useful microbial context but has a different engineering boundary from municipal activated sludge. \citet{VanKessel2015} demonstrate complete nitrification within one microorganism. This evidence limits a literal biological interpretation of a model that assigns ammonium and nitrite oxidation to separate populations. The model can represent two functional conversions without claiming that all real communities organize those conversions in the same way. A learned population coefficient inherits the model's interpretation rather than providing an independent identification of microbial ecology.

Anaerobic ammonium oxidation illustrates a further limit. \citet{Strous1998} investigate slowly growing microorganisms that perform this conversion. \citet{Dietl2015} examine hydrazine synthesis, and \citet{Maalcke2016} characterize an enzyme involved in nitrogen-gas production from hydrazine. These studies concern a distinct metabolic route. A common final nitrogen product does not make its intermediate states and kinetics interchangeable with heterotrophic denitrification. Similarly, the thermophilic treatment studied by \citet{Vandekerckhove2018}, the nitrogen-cycling strategies of \citet{Xie2021}, and the removal and recovery directions reviewed by \citet{WinklerStraka2019} address different conditions and purposes. Mass conservation and non-negativity do not give a model these unrepresented pathways.

Reaction representation alone is also insufficient when transport limits activity. \citet{StenstromSong1991} distinguish oxygen-transport limitation from biological nitrification capacity. \citet{Yilmaz2007} examine biomass activity under alternating conditions during starvation. Their work makes current oxygen availability and prior exposure different sources of information. A steady-state response conditioned on current inputs cannot reproduce every transient effect of that history. A comparison with dynamic control should therefore identify which memory and time-dependent behavior the predictor is able to represent.

Attached-growth treatment gives a useful contrast with completely mixed suspended growth. \citet{Wik2003} reviews trickling filters and biofilm modeling, while \citet{Zhang1994} examine biofilm density, porosity, and pore structure. These physical characteristics create transport and biomass-retention relationships that one bulk concentration vector does not explicitly resolve. \citet{SadinoRiquelme2023} review integration of biological kinetics with computational fluid dynamics. Their discussion makes the trade-off clear. Spatial detail can address transport questions but adds states and calibration requirements. Evidence from a spatially resolved model cannot be transferred automatically to a completely mixed response, even when both describe organic or nitrogen removal.

Treatment configurations make those differences concrete. \citet{Machdar1997} combine anaerobic pretreatment with aerobic post-treatment for developing-country applications. \citet{Tawfik2006} study an anaerobic process followed by a hanging-sponge unit, and \citet{Tawfik2010} examine the effect of sponge volume. The naturally ventilated tower of \citet{Mahmoud2010} and sponge post-treatment of \citet{Mahmoud2011} introduce other retention and aeration arrangements. These studies link process performance to medium and configuration choices. Their inputs do not have the same physical meaning as the volume, recycle, and aeration controls of a completely mixed activated sludge train.

The structure--function comparison of \citet{Hatamoto2018} connects sponge-reactor organization to microbial communities. \citet{Seggelke1999} analyze dynamic nitrification in a low-loaded trickling filter, while \citet{Hellal2021} simulate passive aeration and \citet{Liang2021} model a staged biological filter. These works concern different spatial and temporal states. The model-based upgrade examined by \citet{DiezMontero2019} and the rotating self-aerated reactor studied by \citet{Luan2023} further connect configuration to nutrient behavior. Their relevance is to establish the importance of declared process boundaries. They provide context for transfer questions rather than additional evaluated systems in the activated sludge investigation.

Constituent form matters within those boundaries. \citet{Kasi2011} model dissolved organic nitrogen in a two-stage trickling filter and \citet{Simsek2012} examine its fate. \citet{BressaniRibeiro2021} investigate inorganic-carbon limitation during nitrogen conversion in sponge-bed treatment. These sources show why one nitrogen total cannot identify every growth requirement or removal mechanism. The technology reviews of \citet{Tyagi2021} and \citet{Nasr2022} connect sponge treatment to energy and decentralized applications. The palm-oil effluent application of \citet{Rapi2021} and the contaminant-specific post-treatment of \citet{Mahmoud2018} also keep wastewater composition explicit. A declared activated sludge component set remains a defined scientific scope, rather than a general model of every industrial contaminant or retained medium.

\subsection{Calibration, observability, and numerical solution supply different evidence}

Calibration connects model quantities to observations, but not every observation identifies every parameter. \citet{Petersen2003} critically review experimental designs for activated sludge calibration. \citet{Vanrolleghem1999} use respirometry to estimate parameters and component combinations. Their emphasis on combinations is consequential. An oxygen-use measurement can constrain biological activity while leaving several constituent allocations unresolved. \citet{Machado2014} compare parameter identifiability with influent and operational uncertainty. A small fitting error can therefore coexist with uncertain internal states and nonunique parameter explanations.

Validation studies also use different reference evidence. \citet{Calise2020} validate an activated sludge implementation against benchmark simulation responses. This tests correspondence with a defined reference calculation. In contrast, \citet{SadriMoghaddam2021} characterize influent, calibrate a full-scale treatment model, and compare it with plant measurements. \citet{Meijer2004} connects metabolic phosphorus-removal modeling, consistency of measurements, calibration, and full-scale validation. The latter studies link process models to observed systems, whereas simulator agreement is conditional on a simulator. These evidence types complement one another and should not be presented as equivalent forms of external confirmation.

\citet{Wu2016} simulate and optimize a coking-wastewater process, while \citet{Jasim2020} develops a treatment-plant design model. \citet{Szelag2022} examine nutrient removal and energy under uncertainty, including sensitivity and parameter assessment. \citet{Nazif2023} combine calibrated simulation, influent grouping, and operating optimization. These studies already establish direct model-based routes to engineering decisions. A statistical surrogate can reduce repeated evaluation cost, but it does not create the ability to optimize or remove the uncertainties of the underlying model. \citet{Rivas2008} likewise connect model-based analysis to plant design, where the design question determines which uncertainty matters.

Numerical analysis addresses a separate requirement. \citet{Nelson2009} analyze ASM1 mathematically. \citet{CurtissHirschfelder1952} establish the stiff-equation integration problem, and \citet{KelleyKeyes1998} examine pseudo-transient continuation toward a steady solution. These foundations explain why an apparently plausible parameter set still needs a reliable state calculation. Convergence to a solution, mass conservation of that solution, parameter identification, and agreement with measurements answer different questions. A surrogate trained on numerically accepted reference states can inherit reference-model assumptions and acceptance restrictions even when its statistical fitting is sound.

\section{Prediction tasks determine what a statistical comparison can establish}

\subsection{Soft sensing, component estimation, and prospective prediction}

The review by \citet{Newhart2019} connects wastewater data analysis to data quality, process knowledge, and validation. \citet{Sundui2021} survey machine learning across biological treatment tasks. Their breadth is useful, but it also makes an undifferentiated accuracy ranking misleading. Monitoring an existing condition, estimating an unmeasured constituent, predicting a response at an untested control, and diagnosing a fault use different available information. Their assessment targets cannot be combined into one claim about treatment-model adequacy.

The soft sensor of \citet{Ching2022} illustrates the distinction directly. Its effluent biochemical oxygen demand estimate uses other effluent measurements. Those inputs can be appropriate for estimating a concurrently unmeasured quantity. They do not establish a prospective mapping from fresh influent and selected controls to an effluent that has not yet occurred. \citet{ManavDemir2024} compare learning methods for effluent nutrient-removal parameters. The supervised quantities similarly define the question answered by their error scores. A predictor can be useful within that task without returning a complete component state that satisfies mass conservation and non-negativity.

Sludge and fractionation tasks expose the same issue from another direction. \citet{Ekinci2023} use feature selection and learning for sludge-production estimation, with information that includes observed removal quantities. \citet{Wang2024} predict COD component parameters from a measured COD quantity using separate neural predictors. A production estimate, an influent fractionation estimate, and a future effluent state have different roles in a treatment calculation. Separately accurate fraction estimates also need a consistency check if they are later used together. This follows from the representation and does not require an allegation of improper fitting in the original studies.

The treatment-plant dataset of \citet{Poch1993} was assembled for operational-state classification and fault detection. Measurements from several process stages and removal-performance quantities are therefore meaningful diagnostic inputs. Their availability during diagnosis does not make them available before a new operating choice. Input admissibility must be determined from the prediction time and task. This is a more precise criticism than labeling a measured variable as inherently invalid. It also keeps that observational dataset distinct from model-generated component pairs.

Automation changes the cost of collecting such pairs rather than their scientific meaning. \citet{Caro2024} provide automated acquisition of simulator responses. The contribution supports organization of repeated simulations, while physical acceptance and field agreement remain separate issues. \citet{BaAlawi2021} validate influent sensor information through denoising models. Their contribution concerns the credibility of observed inputs. A simulation surrogate and a sensor model can therefore fail for different reasons even when both are called data-driven treatment models.

Table~\ref{tab:rrl_prediction_evidence} compares these task boundaries. The rows identify useful forms of evidence while keeping their unestablished implications visible. They explain why component-level assessment requires more than importing the best score from an indicator-prediction study.

\begin{table}[htbp]
\centering\footnotesize
\caption{Prediction tasks and the evidence that transfers between them.}
\label{tab:rrl_prediction_evidence}
\begin{tabular}{>{\raggedright\arraybackslash}p{0.23\linewidth}>{\raggedright\arraybackslash}p{0.34\linewidth}>{\raggedright\arraybackslash}p{0.31\linewidth}}
\toprule Task and examples & What the assessment supports & Remaining scope limit\\\midrule
Reference-model comparison\newline\citep{Calise2020,Henze2006} & Agreement with a declared simulated process response. & Reference agreement does not establish field validity or every unrepresented pathway.\\
Calibrated plant modeling\newline\citep{Meijer2004,SadriMoghaddam2021} & Agreement with observations under the calibration and validation conditions. & Several internal states or parameter combinations may remain uncertain.\\
Concurrent soft sensing\newline\citep{Ching2022} & Estimation of an unavailable measurement from information present at that time. & Downstream observations may be unavailable for a prospective operating decision.\\
Component fraction estimation\newline\citep{Wang2024} & Estimation of constituent allocations used in a process model. & Separate component estimates need compatible totals and bounds.\\
Plant-state diagnosis\newline\citep{Poch1993} & Classification of an observed condition and detection of faults. & Diagnostic inputs do not define a forward control-response model.\\
Sludge and effluent indicators\newline\citep{Ekinci2023,ManavDemir2024} & Accuracy for the specified production or quality quantity. & Indicator accuracy leaves full-state mass conservation and non-negativity unresolved.\\
\bottomrule
\end{tabular}
\end{table}

The limitation in the last column is a limit on inference, not a dismissal of the task. A monitoring model can be effective for monitoring while providing little evidence about response at an unobserved control. Using the appropriate comparison target protects both the original contribution and the new engineering claim.

\subsection{Transparent approximations and nonlinear extensions}

Local models and explicit response surfaces give an inspectable starting point for nonlinear treatment behavior. \citet{Li2020} use fuzzy clustering for dissolved-oxygen predictive control, while \citet{Yang2014} organize predictive control around a fuzzy process model. Local representations can distinguish operating regimes, but their transfer depends on coverage of those regimes. A single global predictor has a different approximation burden. Neither choice removes the need to identify the inputs available to the controller.

Polynomial studies address another part of that burden. \citet{Kim2009} examine nitrogen and phosphorus operating optimization. \citet{Lim2011} study organic and nitrogen removal in a membrane bioreactor, and \citet{Ahn2014} compare modeled and experimental operating factors for livestock wastewater. \citet{Nazlabadi2021} review response surfaces in biological wastewater treatment. Explicit curvature and interaction terms can be useful within a sampled domain. Their interpretability does not establish that the selected degree represents a change of reaction regime or an extrapolated condition. These studies therefore support transparent approximation while also defining the domain dependence of that transparency.

Latent-variable regression offers a different compromise. \citet{Wold2001} explain partial least squares and \citet{Khatri2020} review water-quality applications. Shared latent directions can use correlated inputs efficiently, but a linear latent response does not itself represent every nonlinear process relationship. \citet{Woo2009} introduce kernel partial least squares for industrial process variables. \citet{Liu2021} use a dynamic concurrent kernel approach for process monitoring, and \citet{Yang2021} combine latent-variable regression with a relevance-vector model for adaptive effluent prediction. These extensions increase flexibility or accommodate change. They also change the information and time structure of the prediction, which must remain explicit in a steady-state comparison.

Indicator studies remain valuable within that distinction. \citet{Bagherzadeh2021} combine TN prediction with feature selection. \citet{SadriMoghaddamMesghali2023} use an ensemble for TN in a combined trickling-filter and activated sludge system. \citet{Wang2025} compare interpretable learning models with a mechanistic nitrogen predictor. \citet{Yang2025} address waste-sludge generation, while \citet{Xiao2026} connect mechanistic dynamics to interpretable surrogate modeling. Their different targets show how learning can support several treatment decisions. They do not make accuracy for TN, sludge production, or a dynamic indicator interchangeable with full component-state accuracy.

\subsection{Learning families trade flexibility against different forms of structure}

Statistical capacity should be compared rather than assumed from a model's complexity. \citet{Hastie2009} explain regularization and model assessment. \citet{ShwartzZiv2022} question the assumption that deep learning is required for structured tabular data, while \citet{Borisov2024} review the conditions and methods involved in tabular neural learning. Their evidence motivates including simple and flexible families in a common task. It does not justify a universal ordering across wastewater datasets, targets, or training sizes.

Tree ensembles partition the input domain. \citet{Breiman2001} reduces single-tree instability through random forests, and \citet{Geurts2006} increase split randomization in extremely randomized trees. Boosting instead builds the representation sequentially. \citet{Drucker1997} adapts that idea to regression. \citet{ChenGuestrin2016} develop extreme gradient boosting (XGBoost), \citet{Ke2017} develop the light gradient boosting machine (LightGBM), and \citet{Prokhorenkova2018} develop CatBoost with ordered handling of prediction shift. These algorithms offer different regularization and fitting procedures. A tree partition can model nonlinear interactions while leaving coupled output balances unspecified. Its predictive flexibility is therefore a different property from physical enforcement.

Local and kernel learners provide complementary structure. \citet{Mack1981} analyzes nearest-neighbor regression, which ties a prediction to the distribution of nearby observations. \citet{SmolaScholkopf2004} explain support vector regression, where a kernel and regularized fitting objective control response flexibility. The first method makes sparse local coverage visible. The second can represent smooth nonlinear relations without making distance from training support an assurance of accuracy. Both require separate evidence when used beyond the observed input domain.

Multiresponse regression shares information between outputs. \citet{Zou2005} combine coefficient shrinkage and selection through the elastic net. \citet{Yuan2006} select grouped variables, and \citet{Argyriou2008} learn shared features across tasks. The blockwise methods of \citet{Liu2009} and \citet{Simon2013} support grouped multiresponse penalties. These constructions can expose a common statistical structure across components. A shared selected feature or sparsity pattern is not a stoichiometric inventory relation. The useful comparison is whether shared fitting improves the defined response without being mistaken for mass conservation.

Neural regression offers another form of shared flexibility. \citet{Rumelhart1986} develop back propagation and \citet{Hornik1989} establish neural approximation properties under stated conditions. Those representation results are different from a finite-data guarantee about accuracy, optimization, non-negativity, or extrapolation. \citet{Curteanu2011} examine network topology in chemical applications, making architecture selection part of the evidence required for a model comparison. \citet{ArikPfister2021} use sequential feature attention in TabNet. Its masks identify input use within that architecture, whereas named regression terms identify an explicit algebraic structure. Neither form of explanation alone establishes that the returned component state is physically admissible.

\subsection{Data coverage and assessment design limit predictive claims}

Training behavior is not captured by a single final error. \citet{Geman1992} explain the bias--variance dilemma, while \citet{Belkin2019} show that modern learning behavior can depart from a simple classical trade-off. \citet{VieringLoog2023} review the shapes and interpretation of learning curves. These contributions support assessing how performance changes with the amount of labeled information. A low error at one training size cannot establish sample efficiency across sizes, and a learning-curve summary inherits its chosen error scale and assessment design.

Obtaining a label can itself require an expensive reference calculation. \citet{McKay1979} motivate stratified coverage of selected input ranges. Such coverage does not ensure equal representation of every joint biological regime or make independent input intervals representative of a real influent distribution. \citet{Guyon2011} examine active learning when labels are limited, while \citet{Cawley2011} emphasizes regularized baseline comparisons in that setting. These sources identify a separate question about how targets are acquired. Performance with an existing collection of simulator responses does not establish the efficiency of an adaptive sampling strategy.

Model-selection procedures also affect the claim. \citet{Bergstra2011} develop model-based hyperparameter search and \citet{Akiba2019} organize hyperparameter optimization. \citet{Kaufman2012} explain how information leakage allows assessment information to enter fitting decisions. \citet{Pinheiro2025} examine the effects of feature scaling. Their combined lesson is that transformations, model selection, and final assessment must use information in the appropriate role. Scaling that includes assessment observations can change the fitted model, while a model selected against its final assessment set has already used that set as evidence for selection.

This requirement is task-dependent. A current effluent measurement can be a valid soft-sensor input and an unavailable prospective-control input. Leakage concerns arise from the intended information boundary, not from the variable name alone. A common comparison must therefore make both the learning partition and the prediction-time inputs explicit. The critique applies equally to a transparent regression and a flexible neural predictor.

Ranking uncertainty deserves the same care. \citet{Demsar2006} examines statistical comparison across classification datasets, and \citet{Shevchenko2024} examine benchmarking variability in recommender systems. Their applications differ from treatment regression, but both show why one ranking cannot be separated from its independent evidence units and assessment choices. Repeated partitions of one simulated domain are not independent treatment plants. Paired assessment can control some comparison variability while leaving the common reference-model assumptions intact.

Error scales can conceal another kind of variation. \citet{Tofallis2015} discusses relative prediction accuracy, while \citet{Shahbazi2026} address sparse and extreme targets in imbalanced regression. A low residual concentration can have a modest absolute error and a large relative discrepancy. An aggregate total can also hide errors that compensate across components. These concerns justify component and regime-specific assessment alongside a common summary. They do not identify imbalance as the cause of every low-concentration error or make one normalization universally appropriate.

\section{Physical information establishes different kinds of guarantees}

\subsection{Training guidance, output construction, and learned adjustment}

Physics-informed learning covers several mechanisms rather than one guarantee. \citet{Karniadakis2021} review the use of physical knowledge in scientific machine learning, and \citet{Kashinath2021} examine weather and climate applications. Physics can define inputs, restrict a representation, add training penalties, or constrain returned outputs. The useful criticism is therefore specific. A method should be assessed against the relation it actually imposes and the stage at which it imposes that relation. The presence of physical terminology in a model description does not establish every new prediction's mass conservation and non-negativity.

Residual penalties make disagreement with governing equations costly during fitting. \citet{Raissi2019} incorporate differential-equation residuals into neural learning. \citet{Li2024} apply governing-equation information to environmental unit operations, including an activated sludge reactor. These approaches can guide approximation when response observations are limited. A finite penalty also allows its residual to compete with data fitting and other terms. It does not by construction give each returned sample an exact feasibility certificate. This is a property of the objective, rather than a claim that an uninspected implementation performed poorly.

Output construction can establish a stronger but narrower property. \citet{Beucler2021} distinguish residual penalties from reconstructing dependent outputs through analytic constraint layers. Reconstruction can make selected equalities hold by construction. It does not automatically impose a bound on every completed quantity. \citet{Harder2022} make that distinction visible by comparing equality completion with clipping of reconstructed aerosol states. Completion addresses selected mass inventories, whereas clipping addresses non-negativity and can disturb mass conservation. These separate procedures do not establish a joint guarantee. Equality completion may leave negative concentrations, while clipping may change the inventory.

Stoichiometric output mappings enforce structure in another way. \citet{SturmWexler2020} learn material transfers and map them through balanced reaction structure. \citet{SturmWexler2022} use a fixed stoichiometric layer to preserve atom accounting. This ties predicted changes to a mass conservation relation instead of reconstructing selected concentration coordinates. Yet even non-negative transfer magnitudes can consume more of a component than is available. Mass conservation in the component change and non-negativity of the updated state remain separate conditions. A negative tendency can also represent legitimate consumption, so a state bound cannot be applied indiscriminately to every learned quantity.

Learning a constrained adjustment offers a different computational compromise. \citet{Ruckstuhl2021} train a neural map to move unconstrained data-assimilation analyses toward constrained-analysis targets. The targets carry physical information from a more expensive calculation, and a mass penalty can further guide learning. The learned adjustment does not solve the constraint equations for each new analysis. Its reduced violation is therefore distinct from an exact projection at deployment. This comparison identifies where computational savings arise and which guarantee remains to be checked.

Representation-specific construction can sometimes preserve several requirements together. \citet{Geiss2022} impose coarse-to-fine consistency through non-negative multiplicative rescaling, with area weights and a return to dimensional units. The coarse field is preserved, but the constraint does not recover a unique fine-scale distribution or establish an interspecies reaction balance. \citet{Sturm2023} use reduced transport variables and reconstruction to preserve mass conservation and non-negative concentrations for organic aerosol. Both studies show why the accounting operator must match the represented quantity. A spatial average, a transported inventory, and a wastewater constituent state require different weights and boundaries.

Figure~\ref{fig:rrl_enforcement_locations} organizes these distinctions by the location of physical enforcement. The locations can be combined, but they provide different evidence. The figure separates a training preference from an output relation and a checked projection, making the source of a stated guarantee visible.

\begin{figure}[htbp]
\centering
\begingroup\linespread{1}\selectfont
\begin{tikzpicture}[>=Latex,node distance=5mm,
heading/.style={draw=black!75,rounded corners,fill=cyan!6,text width=4.1cm,minimum height=8mm,align=center,font=\footnotesize,inner sep=4pt},
block/.style={heading,fill=white,minimum height=17mm},
limit/.style={block,fill=black!5},
flow/.style={->,thick,draw=black!75}]
\node[heading] (lh) at (-4.55,0) {Construction};
\node[heading] (mh) at (0,0) {Training};
\node[heading] (rh) at (4.55,0) {Returned state};
\node[block,below=of lh] (l1) {Represent process states\\and physical relations};
\node[block,below=of mh] (m1) {Fit observations with\\physical residual penalties};
\node[block,below=of rh] (r1) {Apply a defined\\projection to a\\completed prediction};
\node[block,below=of l1] (l2) {Enforce encoded relations\\within the\\output mapping};
\node[block,below=of m1] (m2) {Penalties guide the fit\\alongside the data objective};
\node[block,below=of r1] (r2) {Impose declared relations\\on each returned state};
\node[limit,below=of l2] (l3) {Guarantee depends on\\the stated representation};
\node[limit,below=of m2] (m3) {Small residuals do not\\certify every new output};
\node[limit,below=of r2] (r3) {Feasibility alone leaves\\process accuracy\\unresolved};
\node[heading,text width=13.3cm,minimum height=12mm,below=7mm of m3] (common)
{Compare mass conservation and non-negativity on the same returned state};
\draw[flow] (lh)--(l1); \draw[flow] (l1)--(l2); \draw[flow] (l2)--(l3);
\draw[flow] (mh)--(m1); \draw[flow] (m1)--(m2); \draw[flow] (m2)--(m3);
\draw[flow] (rh)--(r1); \draw[flow] (r1)--(r2); \draw[flow] (r2)--(r3);
\draw[flow] (l3.south)--(common.north west);
\draw[flow] (m3.south)--(common.north);
\draw[flow] (r3.south)--(common.north east);
\end{tikzpicture}

\endgroup
\caption{Conceptual comparison of where physical knowledge enters a learned model. Construction, fitting, and returned-state projection can be combined. Their guarantees depend on the relations represented and the procedure that establishes them. Mass conservation and non-negativity must be assessed for the same returned state.}
\label{fig:rrl_enforcement_locations}
\end{figure}

The distinction also prevents one mechanism from receiving another mechanism's attribution. Analytic completion, balanced flux mappings, clipping, learned adjustments, and projection have different mathematical roles. Their comparison should preserve those roles before comparing errors or computational cost.

\subsection{State bounds, uncertainty, and approximation capacity}

Non-negativity is a condition on a concentration state. The dynamical-system treatment of \citet{Haddad2010} connects this condition to the direction of the balance at a boundary. A continuous model can preserve it while a numerical step or unconstrained statistical prediction violates it. Clipping changes the returned concentration and may change its inventory. A projection for mass conservation can still return a negative component. The engineering inadequacy is therefore simultaneous rather than separate feasibility. Both requirements must hold after all stages of the prediction procedure.

Strict positivity is a stronger statement because it excludes zero. Zero concentration is allowed by non-negativity and can be an important physical boundary. This difference matters when a relative weighting uses a concentration as a denominator. The weighting may require positive provisional values even though the final physical requirement allows a component to be zero. Naming both requirements precisely avoids treating a numerical safeguard as a physical lower bound.

Probability adds a different object to the assessment. \citet{Hansen2024} condition a predicted mean and covariance on a linear integral mass conservation relation. Exact mass conservation arises in the zero constraint-noise limit, while positive constraint noise allows disagreement. The construction addresses a predictive distribution and a covariance-weighted mean update. It does not by itself establish every component's non-negativity or improve every error metric. \citet{Utkarsh2025} incorporate probabilistic constraint conditioning into learning, distinguishing exact affine distribution updates from approximations used for nonlinear moments. Feasible projected samples and an accurately calibrated nonlinear distribution are different claims.

Ensemble uncertainty has another basis. \citet{Lakshminarayanan2017} use diversity among neural models to estimate predictive uncertainty. Such spread can describe disagreement among predictors without satisfying a material balance. Conversely, a deterministic state that exactly satisfies mass conservation does not provide a calibrated uncertainty distribution. These approaches address complementary inadequacies. Physical constraints restrict what a returned state may contain, while uncertainty assessment concerns how much confidence the data and model support.

Constraint architecture also changes the admissible function class. \citet{Min2024} examine affine and convex enforcement layers and establish neural approximation results under stated feasibility and regularity assumptions. The contribution is stronger than a claim that a penalty often lowers violations, but it applies to the supported neural architecture and constraints. It cannot be transferred to the expressive capacity of a fixed polynomial family or to successful training on finite treatment data. Exact enforcement and adequate representation of the reference response remain separate requirements.

\subsection{Projection depends on its feasible set and projection metric}

Projection uses specified constraints to obtain a returned state from a completed prediction. \citet{SturmSilva2024} use species-weighted least-squares projection for mass conservation through atom accounting. Species scales and error structure affect which component changes are small in that projection metric. The equality constraint still permits negative final concentrations unless a state bound is added. Their work therefore provides a reason to examine the metric, rather than an assurance that every projection for mass conservation is equally accurate or physically complete.

\citet{Valente2025} project predictions onto manifolds defined by selected physical laws, including energy and plasma constraints. This extends the discussion beyond linear mass conservation. A nonlinear feasible manifold need not have a unique nearest point or the global properties of a convex affine set. The local numerical solution and its physical residuals consequently remain part of the evidence. Treating every output projection as the same operation would overlook both the constraint geometry and the numerical problem used to obtain a state.

\citet{KircherVotsmeier2025} combine a projection for mass conservation using relative scaling with interpolation backtracking to address concentration bounds. The provisional values must keep the scaling defined, and the reference must share the required inventory. Interpolation between states satisfying the same mass conservation relation preserves that inventory while controlling the final state bound. When a transferred construction permits zero-valued anchors or boundary contact, its relevant guarantee is non-negativity. The interpolation step is not generally the nearest feasible point under an independently chosen prediction-error metric.

These distinctions limit the interpretation of an accuracy benefit. The convex optimization account of \citet{Boyd2004} gives projection-distance properties for a fixed metric when the reference state belongs to the feasible set. The conditions matter. A prediction-dependent relative weighting followed by interpolation is a different map. An empirical closure may also exclude the mechanistic reference from the constrained set even when that reference satisfies physical identities. Neither situation supports an unconditional claim that projection must improve every component, aggregate metric, or objective.

Projection also imposes only the relations that define its set. A state can conserve the selected inventories and remain non-negative while placing those inventories in the wrong biological forms. The constraints may exclude impossible material creation without selecting the correct kinetic response. This is the central reason to assess mass conservation, non-negativity, and approximation error separately. It identifies a limitation in the scope of the guarantee without denying the value of the guarantee itself.

\section{Interpretability must describe the returned prediction}

\subsection{Named terms and feature use answer different questions}

\citet{Rudin2019} argues for inspecting an interpretable model itself when such a representation is available for a consequential decision. A named term can distinguish the effect of loading, operation, curvature, or interaction. \citet{Brunton2016} use sparse candidate functions to identify dynamical relationships. That contribution links selected functions to scientific explanation, but it concerns a different task from approximating a fixed model's steady response. A fitted response coefficient is not automatically a newly identified kinetic constant or a causal effect of an intervention.

The distinction is sharper when inputs are correlated. A latent variable, a shared sparsity pattern, and an attention mask each describe how a model uses information. The neural feature masks of \citet{ArikPfister2021}, the shared-feature construction of \citet{Argyriou2008}, and the explicit regularized terms discussed by \citet{Hastie2009} offer different kinds of inspection. Their explanatory objects should not be ranked by visibility alone. The appropriate question is whether the explanation describes the component response that the engineer actually receives.

\citet{Shen2023} examine differentiable modeling that combines physical and learned components. Their perspective makes the whole prediction chain relevant to interpretation. An output-coupling relation can transform a driver's coefficient into a different net response. Projection can further change the state and its local sensitivity when a bound becomes active. Explaining only the unadjusted regression therefore leaves a part of the deployed model unexplained. An interpretable raw representation remains useful, but its explanation needs a stated connection to the returned prediction.

\subsection{Constrained regression does not remove fitting and numerical questions}

The published activated sludge surrogate of \citet{Selerio2026} separates an explicit second-order regression from constrained deployment. The fitting structure guides a raw prediction, while checked output procedures establish the stated final-state mass conservation and non-negativity. Its contribution is a transparent component surrogate with a defined physical boundary. Those properties do not by construction extend its single-reactor guarantee to recycle mixing, clarifier storage, or a complete operating decision. Data efficiency, predictive accuracy, and final-state feasibility consequently remain different comparison questions.

Numerical treatment also needs separate attribution. \citet{Razaviyayn2013} analyze block successive minimization under regularity assumptions. Convexity of individual updates does not establish global convexity of a joint fit that includes coupled coefficients and fitted states. \citet{Stellato2018} introduce operator-splitting quadratic optimization and \citet{Stellato2020} develop its mathematical and computational treatment. \citet{HuangfuHall2018} develop a parallel dual revised simplex method. These contributions explain how specified subproblems can be solved. A solver's termination message is still different from independently accepted residuals and bounds, and an efficient solver does not repair an inappropriate feasible set.

The unresolved interpretive issue is therefore linked to the physical issue. A transparent representation is most useful when its terms, units, coupling, and final projection have a coherent meaning. Physical enforcement cannot identify a parameter that the observations leave ambiguous, and a visible coefficient cannot establish that the output obeys its declared balances. Assessing these claims together provides more useful engineering evidence than treating either transparency or feasibility as a complete measure of model credibility.

\section{Connected treatment decisions require more than admissible unit outputs}

\subsection{Hydraulics, separation, and storage connect the unit states}

A biological reactor and a clarifier perform different functions within one material system. \citet{HartelPopel1992} represent secondary clarification and sludge thickening. \citet{Takacs1991} develop a dynamic clarification--thickening model, while \citet{JeppssonDiehl1996} examine the propagation of biological components through the clarifier. Their distinction between bulk solids behavior and component transport matters to a plant surrogate. A scalar solids profile supplies settling and storage information, but it does not by itself define every particulate species' transient distribution. Likewise, plausible outlet concentrations do not determine the stored solids mass.

The benchmark plant of \citet{Alex2008} connects biological stages, clarification, and recirculation. This structure exposes an inadequacy that is absent from an isolated indicator model. A stream leaving one unit is also an input to another unit. Separately predicted unit states can each appear plausible while assigning different component flows to that shared stream. Internal circulation also changes the mixer concentration and reactor flow without creating an external pollutant source. Plant-level mass conservation therefore requires consistent connections as well as local reaction accounting.

Solids age makes storage information decision-relevant. \citet{Leu2012} connect solids retention time to oxygen-transfer behavior, effluent quality, and stability. \citet{Miederer2025} organize wastewater energy-management quantities. These contributions show why equal aeration settings or similar effluent concentrations do not settle resource use. Selected reactor volume affects transfer capacity, and solids retention depends on both stored mass and external solids losses. A quality-only surrogate can therefore leave information needed by an operating objective unresolved, even if its quality prediction is accurate.

The nutrient guidance in \citet{USEPA2009} and component-resolved benchmarking of \citet{Zhang2026} support keeping transformation and separation distinct. A phosphorus total in mixed liquor can differ substantially from its discharged total because particulate phosphorus follows the solids. Soluble nutrient removal is a different question. A nonreactive clarifier does not acquire biological conversion because a predictor reports a lower soluble concentration. The meaningful comparison concerns the component basis, the represented unit processes, and the inventory needed for the intended decision.

\subsection{Network synthesis and fixed-plant operation make different claims}

Network optimization makes routing and treatment connections explicit. \citet{Galan1998} optimize distributed treatment networks and \citet{Galan1999} examine their design and synthesis. \citet{GuelliSouza2011} account for differentiated contaminant regeneration, while \citet{YangNetwork2014} include wastewater regeneration models. These contributions identify a dependence between treatment response and allowable stream destinations. A fixed removal fraction can simplify that dependence, but it cannot represent every concentration-dependent biological or separation response. This is a representation limit that affects the optimization model even when the routing balances are exact.

The application defines the contaminants and decisions. \citet{Gutterres2010} study tannery water reuse, \citet{Hansen2018} minimize petrochemical water and wastewater, and \citet{Xu2015} integrate a yeast-production water system. \citet{PadronPaez2020} address sustainable treatment design through multiple objectives, and \citet{Aboagye2021} connect network design to sustainability assessment. Their common contribution is organized comparison of alternatives under declared priorities. The states and topology choices differ from operating analysis in a fixed activated sludge plant. The network literature provides a conceptual precedent without making those other industries or configurations additional evaluated cases.

Uncertainty and global bounds introduce further distinctions. \citet{Tosarkani2020} develop robust multiobjective network design, while \citet{Ye2019} use hybrid particle-swarm search for planning. \citet{RubioCastro2013} address global optimization of property-based interplant integration. \citet{McCormick1976} provides underestimating constructions that explain how a factorable nonconvex formulation can support bounds. A global certificate belongs to the stated mathematical problem and its valid bounds. It does not establish that an approximate treatment relation reproduces an omitted nonlinear process. A heuristic search and a bound-based global method consequently supply different stopping evidence, even when both produce feasible designs.

\citet{DamalerioRemoval2022} connect phosphorus-removal priorities to treatment and resource objectives. \citet{Selerio2022} develops regression-based design for eutrophication abatement. Their design purpose differs from the implementation-barrier discussion of \citet{DamalerioBarriers2022}. The design studies show why treatment and resource preferences must be stated. Their configuration-level decisions also differ from validating a complete component prediction at fixed plant connections. Evidence for one of these purposes should not be used to imply the other.

\subsection{Operating predictors inherit the targets and time structure of their decisions}

Direct mechanistic optimization is an established comparator rather than a capability introduced by machine learning. \citet{Fang2011} integrate simulation and biological nutrient-removal optimization in a full-scale setting. \citet{Araujo2013} examine optimal operation and the selection of economic controlled variables. The calibrated operating model of \citet{Nazif2023} likewise searches a simulator under influent-dependent conditions. These precedents connect treatment and resource objectives through a process representation. A surrogate changes the response and cost seen by the optimizer, so its contribution should be assessed against this existing capability.

Learned operating objectives can focus on selected quantities. \citet{KusiakWei2013} use neural predictions of airflow and effluent indicators within a multiobjective operating search. \citet{Asadi2017} apply data mining to aeration optimization, and \citet{Mao2024} balance effluent quality and aeration energy through artificial intelligence. These studies establish targeted decision applications. Their predicted objectives do not themselves expose the full stoichiometrically coupled component state. This scope limitation follows from the selected representation and does not justify an unsupported accusation about their data handling.

Dynamic control makes a further distinction. \citet{Sadeghassadi2018} use neural models for setpoint design and predictive control. \citet{Bernardelli2020} place a learning model within real-time predictive control, and \citet{Ren2026} study ammonium-based aeration scheduling under influent disturbances. \citet{Niu2022} combine deep learning and multiobjective dynamic optimization. \citet{Croll2023} review reinforcement learning for wastewater control. These approaches respond to temporal state and action information. A steady-state mapping under fixed controls answers a different question and cannot inherit dynamic disturbance rejection or closed-loop safety from those applications.

Other process objectives reinforce the need to compare like tasks. \citet{Jana2022} combine neural and support-vector methods for detergent-effluent optimization. \citet{Waqas2022} predict membrane permeability in a rotating biological contactor, while \citet{Li2022} combine response surfaces and neural models for membrane fabrication. The process-condition discussion of \citet{Shojaei2021} and anaerobic-digestion review of \citet{Poh2016} examine still other operating settings. These contributions demonstrate useful response approximations. Their objective variables, input availability, and underlying processes determine the evidence they supply. They cannot establish full activated sludge component feasibility merely because a statistical model participates in an optimization.

\subsection{Approximation must be managed where a decision uses it}

The process-surrogate reviews of \citet{BhosekarIerapetritou2018} and \citet{McBride2019} connect model construction to feasibility and optimization. \citet{Durkin2024} use unit surrogates within wastewater resource-recovery optimization, and \citet{He2023} use a Kriging surrogate for papermaking treatment and environmental objectives. Their value arises from repeated evaluation within a larger decision problem. A model that approximates average responses well can nevertheless give a poor local constraint or objective description in the region the optimizer selects.

Multistage and external process applications widen that concern. \citet{Tang2026} develop surrogate-based design of a multistage coal-chemical wastewater network. \citet{Tsochatzidi2025} study surrogate optimization in pharmaceutical systems. \citet{StrausSkogestad2018} organize surrogate construction around self-optimizing variables. These approaches show that the reduced representation can be chosen around a particular decision purpose. Their methods remain conditional on that purpose. A reduced quantity adequate for one objective can omit a component or stored inventory needed by another.

Approximation management addresses this local use explicitly. \citet{Alexandrov1998} develop a trust-region framework for approximation models. \citet{Hameed2026} examine gray-box trust-region filter methods with derivative and feasibility information. \citet{Sahigara2012} compare applicability-domain assessments. A distance-based screen identifies input support, while a large projection identifies disagreement between a raw prediction and imposed relations. These are complementary warning signals. Neither is a calibrated probability of treatment accuracy or a guarantee that the omitted kinetics are satisfied. The stronger convergence claims of adaptive approximation methods require their stated approximation and search conditions.

An optimizer deliberately favors regions with a low predicted objective. It can therefore select a favorable approximation error that contributes little to an average holdout score. Mass conservation and non-negativity remove certain impossible states but leave many inaccurate admissible states. This makes selected-decision evidence an intersection of learning and physical assessment rather than a separate performance label. The reference-model comparison advocated by \citet{Pedrozo2025} addresses that distinction without requiring an assumed optimum of the original engineering problem.

\subsection{Numerical structure and comparison boundaries limit optimality claims}

An optimization problem can be difficult even when part of its response calculation is convex. \citet{Colson2007} review nested decision problems. \citet{Tondel2003} analyze parametric quadratic optimization under a defined coefficient structure. \citet{JainKar2017} examine nonconvex learning problems. Their combined lesson is that an exactly solved inner projection does not make the outer operating search globally convex. Changing active bounds can alter local sensitivities, and control-dependent balance coefficients are different from a fixed-matrix parametric formulation.

The nonlinear numerical methods described by \citet{BuzziFerraris2013}, the sequential programming of \citet{Kraft1988}, and the interior-point filter algorithm of \citet{WachterBiegler2006} support different local search routes. \citet{AudetDennis2006} analyze constrained direct-search convergence, while \citet{Ragonneau2022} develops model-based derivative-free methods. A finite search budget or a few tested directions does not inherit an asymptotic stationarity theorem. Similarly, local solver convergence is different from a global optimality certificate. Stating the stopping evidence protects a decision comparison from a stronger conclusion than its numerical route supplies.

The benchmark study of \citet{Bliek2023} makes evaluation cost and search design central to comparison. A surrogate can have a short search interval while requiring expensive development or later verification. A direct route can have a lower reference objective while using a different local basin. \citet{Zhang2026} reinforce the value of common process boundaries and engineering quantities in wastewater benchmarking. A fair assessment should therefore distinguish the native response used during search from the independently assessed treatment outcome, and should keep the reported computational interval explicit.

The reference calculation has its own scope. Agreement at selected controls supports decision fidelity to that model and its assumptions. It does not automatically establish measured-plant performance or every transient stability property. This limit connects back to calibration and data availability. A numerical decision can be internally well supported while the model's external validity remains an independent scientific question.

\section{Synthesis of the evidence and intersecting knowledge gaps}

\subsection{Established contributions leave linked limits on inference}

The literature supplies substantial capabilities in each part of the problem. Component models define meaningful states and balances. Statistical learners approximate several treatment quantities. Architectural relations and projection can establish specified physical constraints. Interpretable representations make some model terms visible. Model-assisted searches support decisions under declared priorities. The unresolved issue is how these capabilities remain credible when they are used together for the same engineering response.

Table~\ref{tab:rrl_evidence_gap_matrix} summarizes the distinction between an established contribution and the additional evidence needed for an engineering inference. Its rows intersect rather than form independent deficiencies. The representation determines the physical operator and the error target. The prediction-time information determines which outputs can be used in a decision. The returned-state procedure affects both interpretation and the optimizer's local response.

\begin{table}[htbp]
\centering\footnotesize
\caption{Evidence-supported contributions and the linked questions they leave unresolved.}
\label{tab:rrl_evidence_gap_matrix}
\begin{tabular}{>{\raggedright\arraybackslash}p{0.24\linewidth}>{\raggedright\arraybackslash}p{0.31\linewidth}>{\raggedright\arraybackslash}p{0.33\linewidth}}
\toprule Contribution and examples & Evidence established by the contribution & Linked question for engineering use\\\midrule
Process and calibration models\newline\citep{Hauduc2013,Machado2014} & States and parameters have meaning within a specified representation and observational design. & Are the predicted components supported by the available information and the intended boundary?\\
Indicator prediction\newline\citep{Ching2022,Wang2025} & Agreement for a stated quality target under the assessed conditions. & Does the complete response remain accurate when components, regimes, and decision-time inputs are considered together?\\
Physical output enforcement\newline\citep{Beucler2021,KircherVotsmeier2025} & Selected equalities or bounds can be established by the specified output procedure. & Do mass conservation and non-negativity hold simultaneously, and what accuracy or kinetics remain outside that guarantee?\\
Inspectable model structure\newline\citep{Rudin2019,Shen2023} & Terms or information paths can explain a declared representation. & Does that explanation remain valid for the coupled and projected response actually returned?\\
Model-assisted operating search\newline\citep{Durkin2024,Pedrozo2025} & Candidate settings can be selected under a model, objective, and search procedure. & Are the connected states and decision outcomes credible on a common independent reference basis?\\
\bottomrule
\end{tabular}
\end{table}

These comparisons locate limitations in the transfer of evidence between tasks. Each contribution addresses a defined purpose. Its limitation appears when evidence from that purpose is used to justify a broader response or decision. The research need concerns those transfers and their relationships. Mass conservation, learning, and operating optimization each have established foundations, while their joint interpretation remains conditional on the represented task.

\subsection{The dependencies among the gaps}

Figure~\ref{fig:rrl_intersecting_gaps} connects the knowledge gaps through one common prediction problem. Statistical agreement, mass conservation and non-negativity, interpretability, and connected decision evidence all depend on a defined process response and available inputs. Removing any one of these conditions leaves a different uncertainty about the intended engineering use.

\begin{figure}[htbp]
\centering
\begingroup\linespread{1}\selectfont
\begin{tikzpicture}[>=Latex,node distance=7mm,
block/.style={draw=black!75,rounded corners,fill=cyan!6,align=center,font=\footnotesize,inner sep=4pt},
evidence/.style={block,text width=4.1cm,minimum height=20mm,fill=white},
flow/.style={->,thick,draw=black!75}]
\node[block,text width=12.9cm,minimum height=12mm] (scope) at (0,0)
{Defined process response and decision-time information};
\node[evidence] (stat) at (-4.55,-2.25)
{Statistical evidence\\Component accuracy\\Coverage and generalization};
\node[evidence] (phys) at (0,-2.25)
{Physical evidence\\Mass conservation\\Non-negativity};
\node[evidence] (interp) at (4.55,-2.25)
{Interpretive evidence\\Meaning of model terms\\Meaning of returned state};
\node[block,text width=10.1cm,minimum height=18mm,below=8mm of phys] (decision)
{Connected decision evidence\\Agreement at shared streams and stored inventories\\Credibility at selected operating conditions};
\node[block,text width=12.9cm,minimum height=14mm,fill=black!5,below=7mm of decision] (gap)
{Joint knowledge gap\\Evidence for all properties under the same engineering question};
\draw[flow] (scope.south west)--(stat.north);
\draw[flow] (scope.south)--(phys.north);
\draw[flow] (scope.south east)--(interp.north);
\draw[flow] (stat.south)--(decision.north west);
\draw[flow] (phys.south)--(decision.north);
\draw[flow] (interp.south)--(decision.north east);
\draw[flow] (decision)--(gap);
\end{tikzpicture}

\endgroup
\caption{Conceptual relationships among the unresolved evidence requirements. Process representation and data availability condition every strand. Prediction agreement, mass conservation and non-negativity, interpretation, and decision fidelity must be assessed together under a common engineering task. The links express dependencies among scientific questions rather than a prescribed implementation sequence.}
\label{fig:rrl_intersecting_gaps}
\end{figure}

Small prediction error does not establish mass conservation and non-negativity. A response satisfying both requirements can remain inaccurate. A transparent raw model can leave the projection of its prediction unexplained. A connected operating search can select a decision that its native approximation evaluates favorably and the original process model evaluates differently. These are linked possibilities because they concern the same returned response and its use, not unrelated shortcomings of separate methods.

\subsection{Knowledge gaps addressed by the dissertation}

The first gap concerns the meaning and adequacy of the complete predicted response. The reviewed soft-sensing and indicator studies establish selected outputs, whereas engineering interpretation may require their component composition and internal relationships \citep{Ching2022,Ekinci2023,Wang2025}. The dissertation addresses the uncertainty about how statistical agreement across that complete response depends on available information, training support, and operating conditions. The need is for evidence that keeps the prediction target and assessment basis consistent rather than transferring a ranking between different tasks.

The second gap concerns the simultaneous physical validity of that response. The literature distinguishes training penalties, equality construction, state clipping, learned adjustments, and projection. Their guarantees concern different objects and assumptions. Joint mass conservation and non-negativity have already been established for specified constructions \citep{KircherVotsmeier2025,Selerio2026}. The remaining question concerns their relation to approximation quality and intended use within a common activated sludge response. The dissertation addresses whether both properties hold for the returned component state and how far the resulting evidence supports its engineering interpretation. This gap includes the limits of the represented boundary. Satisfaction of selected balances cannot be treated as proof of every biological transformation or unrepresented process.

The third gap concerns the meaning of interpretation after the response is completed. Named coefficients and feature-use explanations can illuminate a statistical model, but coupling and projection can alter the output that supports a decision \citep{Rudin2019,Shen2023}. The dissertation addresses the uncertainty about whether the explanation remains coherent for that returned response. The relevant evidence connects the input definition, learned relation, and final physical state, while preserving the distinction between fitted association and an identified biological mechanism.

The fourth gap concerns transfer from a credible unit response to a credible operating decision. Recirculation, separation, and stored solids link several treatment locations. A locally plausible prediction does not establish consistency at those connections, and a feasible native model does not establish the original process outcome at its selected controls \citep{Alex2008,Pedrozo2025,Zhang2026}. The dissertation addresses the relation among connected-state credibility, declared engineering priorities, approximation error at selected conditions, and computational cost. A decision comparison must remain conditional on its process model and search evidence, with measured-plant validity treated as a separate question.

These gaps meet in one question about engineering trust. The response used to judge an operating choice must be statistically supported, physically defined, interpretable in its returned form, and credible at the connected conditions where the choice is made. Evidence for one property cannot substitute for evidence about the others. The dissertation addresses their joint assessment within the declared activated sludge setting. This boundary preserves a useful and testable contribution without extending it to every treatment technology, microbial pathway, or field application described in the wider literature.
